% JAC1964-20261121, /root/arithmetic_sign.
\subsection*{AR-27. Reduction to curves with an actual canonical base point}

\begin{lemma}[AR-27a]
Let $C/K$ be a smooth projective geometrically integral curve of genus
$g\ge1$ over a number field, with $e\in C(K)$.  For every integer
$n\ge3g-1$ there is a finite flat double cover
$\rho:C'\to C$, with $C'$ smooth projective and geometrically
integral, and a point $e'\in C'(K)$, such that
\[
 \rho^*e=2e',\qquad g(C')=n+1,\qquad
 \omega_{C'}\simeq\mathcal O_{C'}(2ne').                 \tag{AR.27a}
\]
Thus $e'$ itself is a canonical degree-one base point:
$(2g(C')-2)e'\sim K_{C'}$.  No extension of $K$ is needed for
this construction.
\end{lemma}

\begin{proof}
Put $B=\mathcal O_C(ne)\otimes\omega_C^{-1}$.  Then
\[
 \deg(B^2(-e))=2n-4g+3\ge2g+1.
\]
Hence $B^2(-e)$ is very ample.  Choose a section of this line
bundle with geometrically reduced zero divisor avoiding $e$.
Such sections form a nonempty open subset by Bertini on the smooth
curve and by generation at $e$.  This open is defined over $K$
and has a $K$-point because $K$ is infinite.  Multiplication by
the canonical section of $\mathcal O_C(e)$ gives a section
$s\in H^0(C,B^2)$ with a simple zero at $e$ and simple zeros
elsewhere.

Define $C'$ as the double cover with algebra
$\mathcal O_C\oplus B^{-1}$, whose multiplication on
$B^{-1}\otimes B^{-1}$ is given by $s$.  It is finite flat of
degree two.  Locally its equation is $y^2=s(t)$, so it is smooth:
away from the branch divisor it is etale, and at a simple zero the
derivative in $t$ is a unit.  The simple zero at $e$ also proves
geometric integrality.  Indeed a square in the function field of
$C_{\overline K}$ has even valuation at every point, whereas the
local defining rational function of this cover has odd valuation
at $e$.  The quadratic extension remains a field after base change
to $\overline K$.  The unique ramified point above $e$ has residue
field $K$, and $\rho^*e=2e'$.

The tautological section $y$ of $\rho^*B$ vanishes exactly on the
ramification divisor $R$ with multiplicity one.  Thus
$\mathcal O(R)=\rho^*B$.  Riemann--Hurwitz, or directly the local
identity $dt=2y\,dy$ up to a unit at a branch point, gives
\[
 \omega_{C'}=\rho^*\omega_C\otimes\mathcal O(R)
            =\rho^*(\omega_C\otimes B)
            =\rho^*\mathcal O_C(ne)
            =\mathcal O_{C'}(2ne').
\]
Taking degrees gives $2g(C')-2=2n$, completing the proof.
\end{proof}

\begin{proposition}[AR-27b]
The full Fermat height gate of AR-25 can be tested on parametrized
curves whose base point $e$ satisfies
$K_C\sim(2g(C)-2)e$, with that point mapping to the origin of
$E^3$.  This imposes a genuine canonical-point restriction on the
auxiliary curves; the parametrization need not be the normalization
of its image.  The construction preserves the sign and nonvanishing of
the relevant projected modified diagonal.

More precisely, keep the data of AR-25, and apply AR-27a to its
pointed curve $(C,e)$.  Write $K$ here for the field over which
those data are defined.  Write $r=\rho^3:C'^3\to C^3$, let
$\psi'=\psi\circ r$, and put
\[
 \delta'=8\delta,\qquad M'=2M,\qquad
 E_{C'}=\frac1{\delta'}\psi'^*\Pi_*\psi'_*.
\]
Then the exact Chow identities are
\[
 r_*\Delta_{e'}=2\Delta_e,
 \qquad
 E_{C'}\Delta_{e'}=\frac14r^*(E_C\Delta_e).             \tag{AR.27b}
\]
Consequently the normalized raw middle heights satisfy
\[
 h_K(E_{C'}\Delta_{e'},E_{C'}\Delta_{e'})
       =\frac12h_K(E_C\Delta_e,E_C\Delta_e).             \tag{AR.27c}
\]
In particular the AR-25 identity remains exactly
\[
 h_K(z,z)=\frac{\delta'}{D(M')^2}
           h_K(E_{C'}\Delta_{e'},E_{C'}\Delta_{e'}).
                                                               \tag{AR.27d}
\]
If an earlier finite extension was used to construct the AR-25
data, these formulas are over that field; normalization means
that no field-degree factor is introduced.
\end{proposition}

\begin{proof}
The map $f'=f\circ\rho$ sends $e'$ to zero, and
$f'_*[C']=2f_*[C]$.  Hence its Ceresa class is $2Mz$.
Each coordinate map $f_j\circ\rho$ has degree $2\deg(f_j)$,
so $\deg\psi'=8\delta$.

Each of the seven terms of the modified diagonal is a copy of
the underlying curve, with some coordinates equal and the others
fixed at the base point.  On every such term the map $r$ induces
$\rho$, of degree two, and it sends the fixed point to $e$.
Proper pushforward therefore gives the first identity of
(AR.27b) term by term, with the same coefficient two for all
seven terms.  Substitution into the definition of $E_{C'}$ yields
\[
 E_{C'}\Delta_{e'}
   =\frac1{8\delta}r^*\psi^*\Pi_*\psi_*r_*\Delta_{e'}
   =\frac14r^*(E_C\Delta_e).
\]
Height adjointness and $r_*r^*=8\,\mathrm{id}$ now give
\[
 h_K(E_{C'}\Delta_{e'},E_{C'}\Delta_{e'})
 =\frac1{16}h_K(r^*E_C\Delta_e,r^*E_C\Delta_e)
 =\frac8{16}h_K(E_C\Delta_e,E_C\Delta_e).
\]
This proves (AR.27c).  The prefactor changes by the inverse
amount, since $\delta'/(M')^2=2\delta/M^2$, proving (AR.27d).
Also $r^*$ is injective on rational Chow groups, by the same
degree-eight identity.  Thus nonvanishing is preserved.

Every $z$ has an AR-25 presentation, and the preceding construction
turns it into a parametrized presentation with a literal canonical
point.  The map $f\circ\rho$ has degree two onto the previous
image and is not asserted to be its normalization.  Conversely,
the correspondence proof of AR-25 applies to any such parametrized
presentation with $\operatorname{Ce}(f_*[C])=Mz$: its pushforward
identities, finite product-map degree and height adjunction do not
require degree one onto that image.  Thus full-$N$ positivity implies
positivity on this test family, while (AR.27d) gives the reverse
implication.  This does not replace the
projected height by the total height of $\Delta_{e'}$: the
projector remains present in (AR.27c)--(AR.27d).
\end{proof}

\paragraph{Sources and scope.}
The cover is explicitly constructed above.  The curve facts used
are the usual degree-$2g+1$ very-ampleness criterion, Bertini as
in AR-21, and the ramification identity in the Stacks Project,
\S53.12, Tag \texttt{0C1B}, especially Lemmas 53.12.2--53.12.4.
That primary statement and its local differential calculation were
read directly.  Height adjunction is the same input already
checked in AR-24--AR-25.  The new deduction is the canonical-point
cover with its exact Chow and height factors; no positivity
theorem is being added for the projected cycle.
